\chapter{亲核取代}\label{ch:b1:nucleophilic-substitution}

取一瓶 \cip{S}-2-溴丁烷，它是一种单一\omterm{def:b1:stereochemistry-in-depth:enantiomers}{对映异构体}，能使
偏振光旋转。将其与氢氧化钠在极性\omterm{def:b1:intermolecular-forces-solvents:solvent-classes}{非质子性溶剂}中加热，
得到的丁-2-醇同样为单一\omterm{def:b1:stereochemistry-in-depth:enantiomers}{对映异构体} ——但
另一个，即 \cip{R}：反应使分子像被风吹翻的
伞一样内外翻转。改在温水中进行，同一溴代物则得到
部分外消旋的丁-2-醇。同样的总变化，即一个 OH 基团
取代一个 Br 原子，却沿着两条不同的途径进行。本章
由\omterm{def:b1:rate-laws:order}{速率方程}建立这两种\omterm{def:b1:elementary-steps:elementary-step}{反应机理}，阐明它们的
立体化学特征，并给出决定取哪一种机理的规则 ——
这是对动力学与电子效应两章所学工具的首次全面
运用。

\begin{recall}
\omterm{def:b1:rate-laws:order}{速率方程}、\omterm{def:b1:elementary-steps:elementary-step}{基元反应}、\omterm{def:b1:elementary-steps:rds}{决速步}与
\omterm{def:b1:elementary-steps:qssa}{稳态近似}见\cref{ch:b1:rate-laws,ch:b1:elementary-steps}；
\omterm{def:b1:stereochemistry-in-depth:isomers}{构型}、\omterm{def:b1:stereochemistry-in-depth:representations}{楔形式}与\omterm{def:b1:stereochemistry-in-depth:optical-activity}{旋光性}见
\cref{ch:b1:stereochemistry-in-depth}；\omterm{def:b1:electronic-effects:nucleophile}{亲核试剂}、\omterm{def:b1:electronic-effects:nucleophile}{亲电试剂}、
\omterm{def:b1:electronic-effects:nucleophile}{离去基团}和\omterm{def:b1:electronic-effects:intermediates}{碳正离子}见\cref{ch:b1:electronic-effects}；
质子性与\omterm{def:b1:intermolecular-forces-solvents:solvent-classes}{非质子性溶剂}见\cref{ch:b1:intermolecular-forces-solvents}。
\end{recall}

\section{取代反应及其参与物}

\begin{definition}[亲核取代，底物]\label{def:b1:nucleophilic-substitution:substitution}
\emph{亲核取代}\index{亲核取代}是指在饱和碳原子上，
以\omterm{def:b1:electronic-effects:nucleophile}{亲核试剂} Nu 取代\omterm{def:b1:electronic-effects:nucleophile}{离去基团} Y 的过程：
$\mathrm{Nu^-} + \mathrm{R{-}Y} \to \mathrm{R{-}Nu} +
\mathrm{Y^-}$。化合物 R--Y 称为\emph{底物}\index{底物}，
连接 Y 的碳因 C--Y 键向 Y 极化，
而具有亲电性。
\end{definition}

卤代烷是典型\omterm{def:b1:nucleophilic-substitution:substitution}{底物}，以 \ce{Cl-}、\ce{Br-} 或
\ce{I-} 作\omterm{def:b1:electronic-effects:nucleophile}{离去基团}。氢氧根生成醇，醇盐生成
醚；氰根生成腈（多一个碳）；氨和胺生成
胺；碘离子与其他卤离子发生交换。水与醇用作
溶剂时，也能充当\omterm{def:b1:electronic-effects:nucleophile}{亲核试剂}，此时的取代即为
\omterm{def:b1:nucleophilic-substitution:sn1}{溶剂解}。

\section{SN2 机理}

溴甲烷在水中与氢氧根反应时，任一反应物浓度加倍，速率
随之加倍，即 $v = k[\ce{CH3Br}][\ce{OH-}]$，
这是对每种反应物均为一级的反应。

\begin{definition}[SN2 机理，瓦尔登翻转]\label{def:b1:nucleophilic-substitution:sn2}
\emph{SN2 机理}\index{SN2 机理}（取代、亲核、双分子）
是一步\omterm{def:b1:elementary-steps:elementary-step}{基元反应}，其中\omterm{def:b1:electronic-effects:nucleophile}{亲核试剂}
从\omterm{def:b1:electronic-effects:nucleophile}{离去基团}的对侧进攻碳原子，
同时\omterm{def:b1:electronic-effects:nucleophile}{离去基团}离去。碳原子上的另外三个取代基
会像被翻转的伞一样翻到另一侧：这称为\emph{瓦尔登
翻转}\index{瓦尔登翻转}。
\end{definition}

\begin{omfigure}
\setchemfig{atom sep=2.1em}
\schemestart
\chemfig{H@{nu}O^{-}}
\arrow{0}[,0.35]
\chemfig{@{c}C(-[:150]H_3C)(<[:210]H)(<:[:250]C_2H_5)-[@{cb}]@{br}Br}
\arrow{->}
\chemfig{[:0]HO-C(-[:30]CH_3)(<[:-30]H)(<:[:-70]C_2H_5)}
\+
\chemfig{Br^{-}}
\schemestop
\chemmove{\draw[->, thick, shorten <=2pt, shorten >=3pt](nu) .. controls +(15:5mm) and +(175:5mm) .. (c);
          \draw[->, thick, shorten <=2pt](cb) .. controls +(90:6mm) and +(90:6mm) .. (br);}

\medskip
\begin{tikzpicture}[font=\small]
  \node[draw, dashed, rounded corners, inner sep=6pt] at (0,0)
    {$\left[\vcenter{\hbox{\ce{HO}$\cdots$\chemfig{C(-[2]CH_3)(-[:-30]H)(-[:210]C_2H_5)}$\cdots$\ce{Br}}}\right]^{-}$};
  \node[below] at (0,-1.1) {过渡态：碳周围五个基团，其中三个共面};
\end{tikzpicture}

{\small 氢氧根与 \cip{S}-2-溴丁烷发生 SN2 反应。\omterm{def:b1:electronic-effects:nucleophile}{亲核试剂}
从溴的对侧进攻碳原子（\omterm{def:b1:electronic-effects:curly-arrow}{弯箭头}）；在\omterm{def:b1:elementary-steps:energy-profile}{过渡态}中，
O--C 键正在形成，C--Br 键正在断裂；
产物为 \cip{R}-丁-2-醇：其余三个基团已翻转到
另一侧。}
\end{omfigure}

\begin{proposition}[SN2 的速率方程]\label{prop:b1:nucleophilic-substitution:sn2-law}
对于 SN2 反应，$v = k[\mathrm{RY}][\mathrm{Nu}]$。
\end{proposition}

\begin{proof}
该机理是单步双分子\omterm{def:b1:elementary-steps:elementary-step}{基元反应}，由\omterm{def:b1:elementary-steps:elementary-step}{基元反应}的质量作用定律
（见\cref{ch:b1:elementary-steps}），其速率
与两反应物浓度之积成正比。
\end{proof}

\begin{proposition}[SN2 的立体化学]\label{prop:b1:nucleophilic-substitution:sn2-inversion}
在\omterm{def:b1:stereochemistry-in-depth:stereocentre}{立体中心}碳上发生的 SN2 反应只得到一种\omterm{def:b1:stereochemistry-in-depth:isomers}{立体异构体}，
其中三个未变化基团的排列发生翻转：这类反应
称为\omterm{def:b1:nucleophilic-substitution:stereospecific}{立体专一反应}。
\end{proposition}

\begin{proof}
\omterm{def:b1:electronic-effects:nucleophile}{亲核试剂}只能接近碳原子的空侧，即与
\omterm{def:b1:electronic-effects:nucleophile}{离去基团}相对的一侧；其余三个基团在\omterm{def:b1:elementary-steps:energy-profile}{过渡态}中经过一个平面，
并翻到另一侧。新键指向
旧键所不在的方向：每个\omterm{def:b1:nucleophilic-substitution:substitution}{底物}分子都得到
与其镜像对称的基团排列。描述符通常会由 \cip{S}
变为 \cip{R}，或相反，但并非总是如此，因为它取决于
\omterm{def:b1:electronic-effects:nucleophile}{亲核试剂}与\omterm{def:b1:electronic-effects:nucleophile}{离去基团}的优先级。
\end{proof}

\begin{definition}[立体专一反应]\label{def:b1:nucleophilic-substitution:stereospecific}
若立体异构的\omterm{def:b1:nucleophilic-substitution:substitution}{底物}各自给出立体异构不同的产物，
即每个\omterm{def:b1:nucleophilic-substitution:substitution}{底物}各给出自己的产物，则称该反应为
\emph{立体专一反应}\index{立体专一反应}。
\end{definition}

\section{SN1 机理}

2-溴-2-甲基丙烷与水反应，其速率不依赖于\omterm{def:b1:electronic-effects:nucleophile}{亲核试剂}的浓度，
即 $v = k[(\ce{CH3})_3\ce{CBr}]$。
因此决定速率的必定是一个不涉及\omterm{def:b1:electronic-effects:nucleophile}{亲核试剂}的步骤。

\begin{definition}[SN1 机理、外消旋化、溶剂解]\label{def:b1:nucleophilic-substitution:sn1}
\emph{SN1 机理}\index{SN1 机理}（取代、亲核、单分子）
包含两步：先是 C--Y 键的慢速\omterm{def:b1:electronic-effects:cleavage}{异裂}，生成
\omterm{def:b1:electronic-effects:intermediates}{碳正离子}，随后\omterm{def:b1:electronic-effects:nucleophile}{亲核试剂}快速加成到\omterm{def:b1:electronic-effects:intermediates}{碳正离子}上。
把单一\omterm{def:b1:stereochemistry-in-depth:enantiomers}{对映异构体}转化为两者混合物的反应，称为
\emph{外消旋化}\index{外消旋化}；
以溶剂作为\omterm{def:b1:electronic-effects:nucleophile}{亲核试剂}的取代反应，称为
\emph{溶剂解}\index{溶剂解}。
\end{definition}

\begin{omfigure}
\setchemfig{atom sep=2.1em}
\schemestart
\chemfig{C(-[2]CH_3)(-[4]H_3C)(-[6]CH_3)-[@{cb2}]@{br2}Br}
\arrow{->[慢]}
\chemfig{C^{+}(-[2]CH_3)(-[:210]H_3C)(-[:-30]CH_3)}
\+
\chemfig{Br^{-}}
\schemestop

\smallskip
\schemestart
\chemfig{C^{+}(-[2]CH_3)(-[:210]H_3C)(-[:-30]CH_3)}
\arrow{->[快][\ce{H2O}]}
\chemfig{C(-[2]CH_3)(-[4]H_3C)(-[6]CH_3)-OH_2^{+}}
\arrow{->[\ce{-H+}]}
\chemfig{C(-[2]CH_3)(-[4]H_3C)(-[6]CH_3)-OH}
\schemestop
\chemmove{\draw[->, thick, shorten <=2pt](cb2) .. controls +(90:6mm) and +(90:6mm) .. (br2);}

\medskip
\begin{tikzpicture}[font=\small]
  \fill[omProp!15, draw=omProp] (0,0.2) ellipse (0.2 and 0.85);
  \fill[omProp!15, draw=omProp] (0,-0.2) ellipse (0.2 and 0.85);
  \draw[thick] (0,0) -- (-1.1,-0.2) node[left] {\ce{CH3}};
  \draw[thick] (0,0) -- (0.95,-0.45) node[right] {\ce{C2H5}};
  \draw[thick] (0,0) -- (0.4,0.5) node[above right] {H};
  \node[fill=white, inner sep=1pt] at (0,0) {\ce{C+}};
  \draw[->, thick, omDef] (-0.9,1.6) node[left] {\ce{H2O}} -- (-0.15,0.95);
  \draw[->, thick, omDef] (0.9,-1.6) node[right] {\ce{H2O}} -- (0.15,-0.95);
  \node[align=center] at (4.6,0) {从任一面进攻：\\ 两种构型都生成，\\ 若两面之间无其他差别，\\ 则为外消旋混合物};
\end{tikzpicture}

{\small 2-溴-2-甲基丙烷在水中的 SN1 \omterm{def:b1:nucleophilic-substitution:sn1}{溶剂解}：慢速电离
生成叔\omterm{def:b1:electronic-effects:intermediates}{碳正离子}，随后被水捕获并失去一个质子。
另一例：由 2-溴丁烷生成的\omterm{def:b1:stereochemistry-in-depth:stereocentre}{手性}仲\omterm{def:b1:electronic-effects:intermediates}{碳正离子}为平面结构，
水从任一面加成的概率相等。}
\end{omfigure}

\begin{proposition}[SN1 的速率方程]\label{prop:b1:nucleophilic-substitution:sn1-law}
对于 SN1 机理，若\omterm{def:b1:electronic-effects:intermediates}{碳正离子}被捕获的速度远快于
它回到\omterm{def:b1:nucleophilic-substitution:substitution}{底物}的速度，则 $v = k_1[\mathrm{RY}]$，
且与\omterm{def:b1:electronic-effects:nucleophile}{亲核试剂}无关。
\end{proposition}

\begin{proof}
步骤：$\mathrm{RY} \rightleftharpoons \mathrm{R^+} + \mathrm{Y^-}$（$k_1$，
$k_{-1}$），$\mathrm{R^+} + \mathrm{Nu} \to \mathrm{RNu}$（$k_2$）。
\omterm{def:b1:electronic-effects:intermediates}{碳正离子}是一种\omterm{def:b1:electronic-effects:intermediates}{活性中间体}，由\omterm{def:b1:elementary-steps:qssa}{稳态近似}
（见\cref{ch:b1:elementary-steps}）可得 $k_1[\mathrm{RY}] =
k_{-1}[\mathrm{R^+}][\mathrm{Y^-}] + k_2[\mathrm{R^+}][\mathrm{Nu}]$，于是
$v = k_2[\mathrm{R^+}][\mathrm{Nu}] = \dfrac{k_1k_2[\mathrm{RY}][\mathrm{Nu}]}
{k_{-1}[\mathrm{Y^-}] + k_2[\mathrm{Nu}]}$。当 $k_2[\mathrm{Nu}] \gg
k_{-1}[\mathrm{Y^-}]$ 时，上式化为 $k_1[\mathrm{RY}]$，即第一步为
\omterm{def:b1:elementary-steps:rds}{决速步}。
\end{proof}

\begin{proposition}[SN1 的立体化学]\label{prop:b1:nucleophilic-substitution:sn1-stereo}
SN1 反应发生在\omterm{def:b1:stereochemistry-in-depth:stereocentre}{手性}碳原子上时，同时生成两种\omterm{def:b1:stereochemistry-in-depth:isomers}{构型}：
该反应不具有立体专一性，并在
\omterm{def:b1:electronic-effects:intermediates}{碳正离子}两面无差别时导致\omterm{def:b1:nucleophilic-substitution:sn1}{外消旋化}。
\end{proposition}

\begin{proof}
\omterm{def:b1:electronic-effects:intermediates}{碳正离子}为平面结构（见\cref{ch:b1:electronic-effects}），其两个面
互为镜像。若\omterm{def:b1:electronic-effects:nucleophile}{离去基团}在捕获之前已远离，
则\omterm{def:b1:electronic-effects:nucleophile}{亲核试剂}以相同速率从任一面加成，得到
两种等量的\omterm{def:b1:stereochemistry-in-depth:enantiomers}{对映异构体}。实际上，离去的离子仍会
在一段时间内屏蔽其中一个面，因而常观察到
略多的\omterm{def:b1:nucleophilic-substitution:sn2}{瓦尔登翻转}产物。
\end{proof}

\begin{omfigure}
\begin{tikzpicture}
\begin{axis}[
  width=0.8\linewidth, height=5.0cm, xmin=0, xmax=10, ymin=-30, ymax=110,
  axis x line=bottom, axis y line=left, xlabel={反应坐标},
  ylabel={能量}, xtick=\empty, ytick=\empty,
  label style={font=\small}, legend style={font=\scriptsize, draw=none, at={(0.98,0.98)}, anchor=north east}]
\addplot[omDef, very thick] table[x=x, y=E] {figdata/out/bachelor-1-energy-profiles-sn2.dat};
\addplot[omProp, very thick, dashed] table[x=x, y=E] {figdata/out/bachelor-1-energy-profiles-sn1.dat};
\legend{SN2：一步, SN1：两步}
\node[font=\scriptsize, anchor=north] at (axis cs:5,68) {碳正离子};
\node[font=\scriptsize, anchor=north] at (axis cs:0.6,-3) {底物 RY};
\node[font=\scriptsize, anchor=south] at (axis cs:9.6,-18) {产物 RNu};
\end{axis}
\end{tikzpicture}

{\small \omterm{def:b1:elementary-steps:energy-profile}{能量曲线}示意图。SN2：一个\omterm{def:b1:elementary-steps:energy-profile}{过渡态}，碳周围五个
基团。SN1：第一步势垒较高（电离，即\omterm{def:b1:elementary-steps:rds}{决速步}），
\omterm{def:b1:electronic-effects:intermediates}{碳正离子}位于较浅的势阱中，随后经过一个较小的势垒，
被其捕获。}
\end{omfigure}

\section{决定机理的因素}

\begin{proposition}[底物]\label{prop:b1:nucleophilic-substitution:substrate}
甲基和伯\omterm{def:b1:nucleophilic-substitution:substitution}{底物}按 SN2 反应；叔\omterm{def:b1:nucleophilic-substitution:substitution}{底物}按 SN1 反应；
仲\omterm{def:b1:nucleophilic-substitution:substitution}{底物}两种机理都可能，取决于其他因素。
烯丙基和苄基\omterm{def:b1:nucleophilic-substitution:substitution}{底物}两种机理都很快。
\end{proposition}

\begin{proof}
SN2 需要碳的背面有空间：每个烷基取代氢都会挤压\omterm{def:b1:elementary-steps:energy-profile}{过渡态}，
\omterm{def:b1:elementary-steps:energy-profile}{过渡态}带有五个基团，使反应变慢；三个基团则完全阻断反应。SN1 需要
稳定的\omterm{def:b1:electronic-effects:intermediates}{碳正离子}：\omterm{def:b1:electronic-effects:intermediates}{碳正离子}的稳定性顺序
（\cref{prop:b1:electronic-effects:carbocation-order}）使电离
对叔\omterm{def:b1:electronic-effects:intermediates}{碳正离子}很快，对仲\omterm{def:b1:electronic-effects:intermediates}{碳正离子}可以发生，对伯和甲基\omterm{def:b1:electronic-effects:intermediates}{碳正离子}则太慢。
烯丙基和苄基\omterm{def:b1:electronic-effects:intermediates}{碳正离子}可通过共振而稳定，
同一 $\pi$ 体系也稳定 SN2
\omterm{def:b1:elementary-steps:energy-profile}{过渡态}。
\end{proof}

\begin{proposition}[离去基团]\label{prop:b1:nucleophilic-substitution:leaving-group}
两种机理中，\omterm{def:b1:electronic-effects:nucleophile}{离去基团}越好（即其生成的碱越弱），反应越快：
$\ce{I-} > \ce{Br-} > \ce{Cl-} \gg \ce{F-}$；
\ce{OH-} 和 \ce{RO-} 不能离去。
\end{proposition}

\begin{proof}
两种机理中，C--Y 键都在\omterm{def:b1:elementary-steps:rds}{决速步}中或之前断裂，
\omterm{def:b1:elementary-steps:energy-profile}{过渡态}在 Y 上带有部分负电荷。
能很好地容纳负电荷的基团，即\omterm{def:b1:predominant-reaction:strong-acid}{强酸}的共轭碱，
会降低该\omterm{def:b1:elementary-steps:energy-profile}{过渡态}的能量。\ce{HI}、\ce{HBr} 和
\ce{HCl} 是\omterm{def:b1:predominant-reaction:strong-acid}{强酸}，\ce{HF} 是\omterm{def:b1:predominant-reaction:strong-acid}{弱酸}（$\mathrm{p}K_a$ 3.2），水和醇
则是极弱的酸；\ce{OH-} 和 \ce{RO-} 是\omterm{def:b1:predominant-reaction:strong-acid}{强碱}，
是较差的\omterm{def:b1:electronic-effects:nucleophile}{离去基团}。
\end{proof}

\begin{proposition}[溶剂]\label{prop:b1:nucleophilic-substitution:solvent}
极性\omterm{def:b1:intermolecular-forces-solvents:solvent-classes}{质子性溶剂}（水、醇）有利于 SN1；极性\omterm{def:b1:intermolecular-forces-solvents:solvent-classes}{非质子性溶剂}
（丙酮、二甲亚砜、二甲基甲酰胺）有利于阴离子
\omterm{def:b1:electronic-effects:nucleophile}{亲核试剂}参与的 SN2。
\end{proposition}

\begin{proof}
SN1 由中性\omterm{def:b1:nucleophilic-substitution:substitution}{底物}生成两个离子：极性\omterm{def:b1:intermolecular-forces-solvents:solvent-classes}{质子性溶剂}
同时稳定两者，阳离子通过其\omterm{def:b1:lewis-resonance-vsepr:lewis}{孤电子对}，阴离子通过\omterm{def:b1:intermolecular-forces-solvents:hbond}{氢键}，
并降低电离势垒。在 SN2 中，阴离子\omterm{def:b1:electronic-effects:nucleophile}{亲核试剂}的电荷在\omterm{def:b1:elementary-steps:energy-profile}{过渡态}中分散；
\omterm{def:b1:intermolecular-forces-solvents:solvent-classes}{质子性溶剂}通过\omterm{def:b1:intermolecular-forces-solvents:hbond}{氢键}强烈结合小的阴离子，
使其反应变慢，
而\omterm{def:b1:intermolecular-forces-solvents:solvent-classes}{非质子性溶剂}使阴离子自由而具有反应活性
（\cref{ch:b1:intermolecular-forces-solvents}）。
\end{proof}

\begin{method}[选择机理]\label{met:b1:nucleophilic-substitution:choose-mechanism}
\begin{enumerate}
  \item \omterm{def:b1:nucleophilic-substitution:substitution}{底物}：甲基或伯 $\to$ SN2；叔 $\to$ SN1；
        仲 $\to$ 继续判断。
  \item \omterm{def:b1:electronic-effects:nucleophile}{亲核试剂}：强且为阴离子（$\ce{OH-}$、$\ce{RO-}$、$\ce{CN-}$、
        $\ce{I-}$）$\to$ SN2；弱且为中性（水、醇）$\to$ SN1。
  \item 溶剂：非质子性 $\to$ SN2；质子性 $\to$ SN1。
  \item 预测立体化学（翻转或\omterm{def:b1:nucleophilic-substitution:sn1}{外消旋化}），
        并检查\omterm{def:b1:predominant-reaction:strong-acid}{强碱}是否更可能引起消除
        （\cref{ch:b1:elimination}）。
\end{enumerate}
\end{method}

\begin{omfigure}
\begin{tabular}{lll}
\hline
 & SN2 & SN1 \\
\hline
步骤 & 一步 & 两步，经由\omterm{def:b1:electronic-effects:intermediates}{碳正离子} \\
\omterm{def:b1:rate-laws:order}{速率方程} & $k[\mathrm{RY}][\mathrm{Nu}]$ & $k[\mathrm{RY}]$ \\
\omterm{def:b1:nucleophilic-substitution:substitution}{底物} & 甲基 $>$ 伯 $>$ 仲 & 叔 $>$ 仲 \\
\omterm{def:b1:electronic-effects:nucleophile}{亲核试剂} & 强，阴离子 & 弱，中性（常为溶剂） \\
溶剂 & 极性非质子性 & 极性质子性 \\
立体化学 & 翻转（立体专一） & \omterm{def:b1:nucleophilic-substitution:sn1}{外消旋化}（大多数情况） \\
\hline
\end{tabular}

\smallskip
{\small 两种机理并列对比。}
\end{omfigure}

\section{练习}

\begin{exercise}[$\star$]\label{exo:b1:nucleophilic-substitution:1}
写出下列反应的产物：溴乙烷与氢氧化钠；碘甲烷与乙醇钠；
1-氯丙烷与氰化钾；溴甲烷与氨（第一步）。
指出\omterm{def:b1:nucleophilic-substitution:substitution}{底物}、\omterm{def:b1:electronic-effects:nucleophile}{亲核试剂}和\omterm{def:b1:electronic-effects:nucleophile}{离去基团}。
\end{exercise}

\begin{exercise}[$\star$]\label{exo:b1:nucleophilic-substitution:2}
对于反应 $\mathrm{RY} + \mathrm{Nu}$，\omterm{def:b1:rate-laws:order}{速率方程}为 $v = k[\mathrm{RY}][\mathrm{Nu}]$，
若 $[\mathrm{Nu}]$ 增至三倍，\omterm{def:b1:rate-laws:initial-rate}{初始速率}如何变化？若两者浓度都减半呢？
若\omterm{def:b1:rate-laws:order}{速率方程}为 $v = k[\mathrm{RY}]$，对同样的问题作答。
\end{exercise}

\begin{exercise}[$\star$]\label{exo:b1:nucleophilic-substitution:3}
预测下列反应的机理（SN1 或 SN2）：1-溴丁烷与碘化钠在丙酮中；
2-溴-2-甲基丙烷在水中；溴甲烷与氢氧化物；
2-溴丁烷在无添加碱的乙醇中。
\end{exercise}

\begin{exercise}[$\star$]\label{exo:b1:nucleophilic-substitution:4}
画出氰根与 \cip{R}-2-溴丁烷发生 SN2 反应的产物的
\omterm{def:b1:stereochemistry-in-depth:representations}{楔形式}（克拉姆表示），并给出其\omterm{def:b1:stereochemistry-in-depth:isomers}{构型}描述符。
\end{exercise}

\begin{exercise}[$\star\star$]\label{exo:b1:nucleophilic-substitution:5}
写出 2-溴-2-甲基丙烷水解的完整 SN1 机理，
用\omterm{def:b1:electronic-effects:curly-arrow}{弯箭头}表示，包括质子的脱去。为何\omterm{def:b1:rate-laws:order}{速率方程}
为一级，尽管有三种物种参与？
\end{exercise}

\begin{exercise}[$\star\star$]\label{exo:b1:nucleophilic-substitution:6}
解释为何 1-溴-2,2-二甲基丙烷（一种伯溴化物）的 SN2 \omterm{def:b1:rate-laws:rate}{反应
速率}极慢，且它不发生 SN1 反应。
\end{exercise}

\begin{exercise}[$\star\star$]\label{exo:b1:nucleophilic-substitution:7}
\cip{S}-2-碘辛烷在丙酮中与碘化钠放置时失去\omterm{def:b1:stereochemistry-in-depth:optical-activity}{旋光性}，
尽管未生成新化合物。解释之，并说明
为何活性丧失速率恰为取代速率的两倍。
\end{exercise}

\begin{exercise}[$\star\star$]\label{exo:b1:nucleophilic-substitution:8}
按与碘离子发生 SN2 反应的速率排序：2-溴-2-甲基丙烷、溴甲烷、
2-溴丙烷、1-溴丙烷。将同一组化合物按 SN1
\omterm{def:b1:nucleophilic-substitution:sn1}{溶剂解}速率排序。
\end{exercise}

\begin{exercise}[$\star\star$]\label{exo:b1:nucleophilic-substitution:9}
一种旋光纯的仲溴化物发生\omterm{def:b1:nucleophilic-substitution:sn1}{溶剂解}，产物
的旋光度为纯翻转所预期数值的 \qty{12}{\percent}，
求出产物两种\omterm{def:b1:stereochemistry-in-depth:enantiomers}{对映异构体}的比例，
并加以解释。
\end{exercise}

\begin{exercise}[$\star\star\star$]\label{exo:b1:nucleophilic-substitution:10}
用\omterm{def:b1:elementary-steps:qssa}{稳态近似}推导一般的 SN1 \omterm{def:b1:rate-laws:order}{速率方程}，
并证明在反应开始时加入\omterm{def:b1:electronic-effects:nucleophile}{离去基团}离子 \ce{Y-} 会减慢
反应（动力学上的\omterm{def:b1:precipitation:common-ion}{同离子效应}）。在何种条件下
此效应可以忽略？
\end{exercise}

\begin{exercise}[$\star\star\star$]\label{exo:b1:nucleophilic-substitution:11}
溴甲烷与氢氧根反应的两个\omterm{def:b1:rate-laws:half-life}{半衰期}：取
$[\ce{OH-}]_0 = [\ce{CH3Br}]_0 = C_0$，证明 $1/[\ce{CH3Br}] = 1/C_0 +
kt$，并给出 $t_{1/2}$。当 $[\ce{OH-}]_0 = 50\,C_0$ 时，证明
衰减为指数型，并给出 $t_{1/2}$。本题数据：$k =
\qty{2.0e-4}{L/(mol.s)}$，$C_0 = \qty{0.010}{mol/L}$。
\end{exercise}

\begin{exercise}[$\star\star\star$]\label{exo:b1:nucleophilic-substitution:12}
解释为何醇 \ce{ROH} 与溴化钠不反应，但
能与浓氢溴酸反应生成 \ce{RBr}。写出
叔醇与伯醇的\omterm{def:b1:elementary-steps:elementary-step}{反应机理}。
\end{exercise}

\section{问题：叔丁基氯的溶剂解}

\begin{problem}[{周末问题——由电导率数据求级数与速率常数、
半衰期、机理及其立体化学，
以及溶剂解的活化能}]
\label{pb:b1:nucleophilic-substitution:1}
将 2-氯-2-甲基丙烷（叔丁基氯）以低浓度溶于
水–丙酮混合物中。其\omterm{def:b1:nucleophilic-substitution:sn1}{溶剂解}
\ce{(CH3)3CCl + H2O -> (CH3)3COH + H+ + Cl-} 生成离子，而
溶液的电导率 $\kappa$，它与所生成离子的量成正比，
被记录。经很长时间后，它
在两个温度下均达到 $\kappa_\infty = \qty{2.000}{mS/cm}$。测量值
（\unit{mS/cm}）：

\begin{center}
\begin{tabular}{lcccccccc}
\hline
$t$（分钟） & 0 & 5 & 10 & 15 & 20 & 30 & 45 & 60 \\
\hline
$\kappa$, \qty{25.0}{\celsius} & 0 & 0.172 & 0.329 & 0.473 & 0.605 & 0.835 & 1.110 & 1.321 \\
$\kappa$, \qty{35.0}{\celsius} & 0 & 0.507 & 0.885 & 1.168 & 1.379 & 1.654 & 1.856 & 1.940 \\
\hline
\end{tabular}
\end{center}
% data generated in figdata/bachelor-1/solvolysis.py (story data)

\medskip
\noindent\textbf{第一部分——级数。}
\begin{enumerate}
  \item 为何电导率与已参与反应的
        \omterm{def:b1:nucleophilic-substitution:substitution}{底物}的量成正比？
  \item 将 $[\mathrm{RCl}]/[\mathrm{RCl}]_0$ 用
        $\kappa$ 与 $\kappa_\infty$ 表示。
  \item 写出一级反应的积分\omterm{def:b1:rate-laws:order}{速率方程}，以及
        为检验之而需对 $t$ 作图的量。
  \item 对 \qty{25.0}{\celsius} 下的每个
        时刻计算 $\ln(\kappa_\infty - \kappa)$。
  \item 所作图是否为直线？由此判断级数。
  \item 水大量过量。这在级数中掩盖了什么？
\end{enumerate}

\medskip
\noindent\textbf{第二部分——\omterm{def:b1:rate-laws:order}{速率常数}与\omterm{def:b1:rate-laws:half-life}{半衰期}。}
\begin{enumerate}[resume]
  \item 求 \qty{25.0}{\celsius} 时的\omterm{def:b1:rate-laws:order}{速率常数}，单位为 \unit{s^{-1}}。
  \item 计算\omterm{def:b1:rate-laws:half-life}{半衰期}。
  \item \omterm{def:b1:nucleophilic-substitution:substitution}{底物}消耗 \qty{90}{\percent} 时是何时刻？
  \item 求 \qty{35.0}{\celsius} 时的\omterm{def:b1:rate-laws:order}{速率常数}。
  \item 温度升高 10 ℃ 时，速率增大到原来的多少倍？
  \item 对照该定律，检验 \qty{35}{\celsius} 系列中的一个数值。
\end{enumerate}

\medskip
\noindent\textbf{第三部分——机理。}
\begin{enumerate}[resume]
  \item \omterm{def:b1:nucleophilic-substitution:substitution}{底物}与溶剂提示何种机理？
  \item 用\omterm{def:b1:electronic-effects:curly-arrow}{弯箭头}写出该机理。
  \item 证明该机理与所得的级数相符。
  \item 加入浓度为 \qty{0.01}{mol/L} 的氢氧化钠，几乎不改变
        氯离子的\omterm{def:b1:rate-laws:rate}{消耗速率}。为何？
  \item 对\omterm{def:b1:stereochemistry-in-depth:stereocentre}{手性}叔氯代烷进行同样的实验，
        旋光纯的 3-氯-3-甲基己烷得到一种几乎
        无\omterm{def:b1:stereochemistry-in-depth:optical-activity}{旋光性}的醇。试解释。
  \item 画出该反应的\omterm{def:b1:elementary-steps:energy-profile}{能量曲线}，并标出
        \omterm{def:b1:elementary-steps:rds}{决速步}。
  \item 在相同条件下，1-氯丁烷为何几乎不反应？
\end{enumerate}

\medskip
\noindent\textbf{第四部分——\omterm{def:b1:rate-laws:activation-energy}{活化能}。}
\begin{enumerate}[resume]
  \item 写出\omterm{thm:b1:rate-laws:arrhenius}{阿伦尼乌斯定律}。
  \item 用两个\omterm{def:b1:rate-laws:order}{速率常数}表示 $E_a$。
  \item 计算 $E_a$（$R = \qty{8.314}{J/(K.mol)}$）。
  \item 在该机理中，这一能量表示什么？
  \item 预测 \qty{45.0}{\celsius} 时的\omterm{def:b1:rate-laws:order}{速率常数}。
  \item 给出\omterm{def:b1:nucleophilic-substitution:sn1}{溶剂解}的\omterm{def:b1:rate-laws:activation-energy}{活化能}，保留两位
        有效数字。
\end{enumerate}
\end{problem}
